%% sections/01_introduction.tex

\section{서론}
\label{sec:introduction}

\subsection{배경과 두 가지 질문}

투기적 디코딩(speculative decoding)~\cite{leviathan2023spec,chen2023spec}은
소형 드래프터가 제안한 복수 토큰을 타깃 모델이 단일 forward pass로 검증하여
autoregressive 생성의 순차성을 깬다.
SuffixDecoding~\cite{snowflake-arctic}, EAGLE~\cite{li2024eagle,li2025eagle3} 등이
프로덕션에 배포되며 spec-decode 는 GPU 기반 LLM 서빙의 \emph{지배적} 처리량 lever 가 되었다.
본 연구의 B200 측정에서도 단일 기법 중 spec-decode 의 이득($+232\%$까지)이
다른 어떤 lever 보다 크다.

그러나 프로덕션 보고들은 이 가속이 워크로드와 동시성에 크게 의존함을
일관되게 드러낸다~\cite{redhat2025eagle3vllm,peagle2026vllm}.
운영자는 이질적 모델 fleet(7B 대화 모델부터 671B MoE 추론 모델까지)과
이질적 트래픽(대화$\cdot$코드$\cdot$추론)을 동시에 서빙한다.
이 현실에서 두 가지 질문이 미해결로 남아 있다.

\begin{itemize}
\item \textbf{Q1 (언제 켜는가).} 주어진 (모델, 워크로드)에서 spec-decode 는
  언제 이득이고 언제 손해인가? 그 부호를 무엇으로 \emph{예측}하는가?
\item \textbf{Q2 (무엇을 남기는가).} spec-decode 의 성공/실패는 하드웨어 자원
  (GPU$\cdot$CPU 이용률)을 어떻게 바꾸며, 그것이 다음 최적화 기회를 어디에 만드는가?
\end{itemize}

\subsection{접근: B200 대규모 특성화}

우리는 DGX B200$\times$8 에서 \textbf{10개 모델}(7B--671B, dense 및 MoE) $\times$
\textbf{4개 방법}(vanilla / suffix / ngram / llm-d) $\times$
\textbf{7개 조건}(실 서빙 trace 6 corpus $+$ 혼합)을 측정하여
Q1$\cdot$Q2 에 데이터로 답한다.
모든 결과는 B200 실측이며, 합성 워크로드를 실제 trace
(ShareGPT, WildChat, LMSYS-Chat, SWE-bench, MBPP, HumanEval)로 대체해
in-distribution 편향을 제거하였다.

\subsection{기여}

\begin{enumerate}
\item \textbf{예측 조건 (Q1).}
  spec-decode 의 부호는 draft 수락률 $\alpha$ 와 per-step overhead 의 경쟁으로
  결정됨을 70개 (모델$\times$조건) 쌍으로 정량화한다.
  net-positive 셀 $\alpha$ 중앙값 $0.68$ vs.\ vanilla-win $0.37$ 이며,
  break-even 은 모델 overhead 에 따라 이동한다(\S\ref{sec:results}).

\item \textbf{실패 메커니즘.}
  vanilla 가 이기는 15개 셀은 \emph{모두} 추론 특화/MoE 모델이다.
  DeepSeek-R1 671B(MoE 추론)는 전 조건 $-40\!\sim\!-64\%$ 손해이며,
  긴 사고 사슬의 출력 다양성 $\times$ MoE expert routing 동적성이
  draft prefix-match 를 무너뜨려 $\alpha$ 를 낮춤을 분석한다
  (\S\ref{sec:mechanism}). 최적 방법은 크기가 아닌 \emph{모델 유형}이 결정한다.

\item \textbf{자원 signature (Q2).}
  spec-decode 가 성공하는 곳에서 GPU 가 un-saturate($94.9\%\!\to\!78.7\%$)되고
  CPU 는 7B--671B 전 구성에서 $\le4.9\%$ 유휴하다.
  처리량 지배 기법이 시스템을 host(CPU)-bound 로 전환함을 보인다.

\item \textbf{CPU 병렬 host-side reclamation (방법 + 성능).}
  위 자원 여유를 회수하기 위해, suffix 가 키운 step당 host-path 작업을
  유휴 CPU 코어에 \emph{병렬} 분산하여 GPU verify 뒤로 숨기는
  reclamation 을 제안하고(\S\ref{sec:direction}, hiding 조건 식~\ref{eq:hide-condition}),
  B200 에서 처리량을 회수한다.
  \todo{회수 처리량 이득(suffix 대비, 기하평균)과 GPU 이용률 회복폭 --- 측정 후 기입
  (\S\ref{subsec:res-reclaim}).}
\end{enumerate}

\subsection{논문 구성}
\S\ref{sec:background} 배경,
\S\ref{sec:related} 관련 연구,
\S\ref{sec:motivation} 특성화 개요,
\S\ref{sec:problem} 예측 조건 정식화,
\S\ref{sec:direction} host-side reclamation 방향,
\S\ref{sec:methodology} 실험 방법,
\S\ref{sec:results} 측정 결과,
\S\ref{sec:mechanism} 메커니즘 분석,
\S\ref{sec:discussion} 논의,
\S\ref{sec:conclusion} 결론.
